% answers.tex —— 2014 年·上海卷·理科 · 参考答案与详细解析片段
% 由 tools/generate_projects.py 依来源总账自 2014/2014_数学真题_上海理_试题.md + 2014/2014_数学真题_上海理_解析正文.md + 2014/2014_数学真题_上海理_答案要点.md 生成，勿手改（改源 MD 后重跑）


\begin{answer}
\gksolhead{1}{}

\textbf{题目}：

函数 $y=1-2\cos^2(2x)$ 的最小正周期是\gkblank{}。

\Needspace{4\baselineskip}
\textbf{解析}

\Needspace{5\baselineskip}
由二倍角公式 $\cos^2(2x)=\dfrac{1+\cos4x}{2}$，得
\nopagebreak[4]
\[y=1-2\cos^2(2x)=1-(1+\cos4x)=-\cos4x,\]
\Needspace{5\baselineskip}
故最小正周期
\nopagebreak[4]
\[T=\frac{2\pi}{4}=\frac\pi2 .\]

\gkanswer{答案：$\dfrac\pi2$。由 $y=1-2\cos^2(2x)=-\cos4x$ 得 $T=\dfrac{2\pi}4=\dfrac\pi2$。}

\end{answer}
\begin{answer}
\gksolhead{2}{}

\textbf{题目}：

若复数 $z=1+2\mathrm i$，其中 $\mathrm i$ 是虚数单位，则 $\left(z+\dfrac1{\overline z}\right)\cdot\overline z=$\gkblank{}。

\Needspace{4\baselineskip}
\textbf{解析}

\Needspace{5\baselineskip}
由 $z=1+2\mathrm i$ 得 $\overline z=1-2\mathrm i$，$|z|^2=1^2+2^2=5$。利用 $\overline z\cdot\dfrac1{\overline z}=1$，
\nopagebreak[4]
\[\left(z+\frac1{\overline z}\right)\overline z=z\overline z+1=|z|^2+1=5+1=6 .\]
也可先算 $\dfrac1{\overline z}=\dfrac{z}{|z|^2}=\dfrac{1+2\mathrm i}5$，再乘 $\overline z$ 得 $\dfrac{(1+2\mathrm i)(1-2\mathrm i)}5=1$，结果同为 $6$。

\gkanswer{答案：$6$。$\left(z+\dfrac1{\overline z}\right)\overline z=|z|^2+1=5+1=6$。}

\end{answer}
\begin{answer}
\gksolhead{3}{}

\textbf{题目}：

若抛物线 $y^2=2px$ 的焦点与椭圆 $\dfrac{x^2}{9}+\dfrac{y^2}{5}=1$ 的右焦点重合，则该抛物线的准线方程为\gkblank{}。

\Needspace{4\baselineskip}
\textbf{解析}

\Needspace{5\baselineskip}
椭圆 $\dfrac{x^2}9+\dfrac{y^2}5=1$ 中 $a^2=9$，$b^2=5$，故
\nopagebreak[4]
\[c=\sqrt{a^2-b^2}=\sqrt{9-5}=2,\]
右焦点为 $(2,0)$。

\Needspace{5\baselineskip}
抛物线 $y^2=2px$ 的焦点为 $\left(\dfrac p2,0\right)$，由 $\dfrac p2=2$ 得 $p=4$，其准线为
\nopagebreak[4]
\[x=-\frac p2=-2 .\]

\gkanswer{答案：$x=-2$。椭圆右焦点为 $(2,0)$，故 $\dfrac p2=2$，$p=4$，准线 $x=-\dfrac p2=-2$。}

\end{answer}
\begin{answer}
\gksolhead{4}{}

\textbf{题目}：

\Needspace{5\baselineskip}
设
\nopagebreak[4]
\[f(x)=\begin{cases}
x, & x\in(-\infty,a),\\
x^2, & x\in[a,+\infty),
\end{cases}\]
若 $f(2)=4$，则 $a$ 的取值范围为\gkblank{}。

\Needspace{4\baselineskip}
\textbf{解析}

若 $a>2$，则 $2\in(-\infty,a)$，此时 $f(2)=2\neq4$，不合题意。

若 $a\leqslant2$，则 $2\in[a,+\infty)$，此时 $f(2)=2^2=4$ 恒成立。

故 $a$ 的取值范围为 $(-\infty,2]$。

\gkanswer{答案：$(-\infty,2]$。$a\leqslant2$ 时 $2\in[a,+\infty)$，$f(2)=2^2=4$；$a>2$ 时 $f(2)=2\neq4$。}

\end{answer}
\begin{answer}
\gksolhead{5}{}

\textbf{题目}：

若实数 $x$，$y$ 满足 $xy=1$，则 $x^2+2y^2$ 的最小值为\gkblank{}。

\Needspace{4\baselineskip}
\textbf{解析}

\Needspace{5\baselineskip}
由 $xy=1$ 知 $x$，$y$ 同号且都不为零，于是
\nopagebreak[4]
\[x^2+2y^2\geqslant2\sqrt{x^2\cdot2y^2}=2\sqrt2\,|xy|=2\sqrt2 .\]
\Needspace{5\baselineskip}
等号成立当且仅当 $x^2=2y^2$ 且 $xy=1$，即 $y^2=\dfrac1{\sqrt2}$，$x=\sqrt2\,y$，例如
\nopagebreak[4]
\[x=2^{\frac14},\qquad y=2^{-\frac14},\]
此时 $xy=1$ 且 $x^2+2y^2=\sqrt2+\sqrt2=2\sqrt2$。故最小值为 $2\sqrt2$。

\gkanswer{答案：$2\sqrt2$。$x^2+2y^2\geqslant2\sqrt2|xy|=2\sqrt2$，$x^2=2y^2$ 且 $xy=1$ 时取等号。}

\end{answer}
\begin{answer}
\gksolhead{6}{}

\textbf{题目}：

若圆锥的侧面积是底面积的 $3$ 倍，则其母线与底面夹角的大小为\gkblank{}（结果用反三角函数值表示）。

\Needspace{4\baselineskip}
\textbf{解析}

\Needspace{5\baselineskip}
设底面半径为 $r$，母线长为 $l$，母线与底面所成角为 $\theta$。圆锥侧面积为 $\pi rl$，底面积为 $\pi r^2$，由题意
\nopagebreak[4]
\[\pi rl=3\pi r^2\iff l=3r .\]
\Needspace{5\baselineskip}
母线与底面夹角满足
\nopagebreak[4]
\[\cos\theta=\frac rl=\frac13,\]
故 $\theta=\arccos\dfrac13$。

\gkanswer{答案：$\arccos\dfrac13$。由 $\pi rl=3\pi r^2$ 得 $l=3r$，故 $\cos\theta=\dfrac rl=\dfrac13$。}

\end{answer}
\begin{answer}
\gksolhead{7}{}

\textbf{题目}：

已知曲线 $C$ 的极坐标方程为 $\rho(3\cos\theta-4\sin\theta)=1$，则 $C$ 与极轴的交点到极点的距离是\gkblank{}。

\Needspace{4\baselineskip}
\textbf{解析}

\Needspace{5\baselineskip}
由 $x=\rho\cos\theta$，$y=\rho\sin\theta$ 得曲线 $C$ 的直角坐标方程
\nopagebreak[4]
\[3x-4y=1 .\]
极轴即射线 $y=0$（$x\geqslant0$）。令 $y=0$ 得 $x=\dfrac13>0$，交点为 $\left(\dfrac13,0\right)$，它到极点的距离为 $\dfrac13$。

\gkanswer{答案：$\dfrac13$。曲线即 $3x-4y=1$，令 $y=0$ 得 $x=\dfrac13$。}

\end{answer}
\begin{answer}
\gksolhead{8}{}

\textbf{题目}：

设无穷等比数列 $\{a_n\}$ 的公比为 $q$，若 $a_1=\lim\limits_{n\to\infty}(a_3+a_4+\cdots+a_n)$，则 $q=$\gkblank{}。

\Needspace{4\baselineskip}
\textbf{解析}

\Needspace{5\baselineskip}
因 $\lim\limits_{n\to\infty}(a_3+a_4+\cdots+a_n)$ 存在，故 $|q|<1$ 且 $q\neq0$；又 $\{a_n\}$ 为等比数列，$a_1\neq0$。由 $a_k=a_1q^{k-1}$ 得
\nopagebreak[4]
\[\lim_{n\to\infty}(a_3+a_4+\cdots+a_n)=\frac{a_3}{1-q}=\frac{a_1q^2}{1-q}.\]
\Needspace{5\baselineskip}
由 $a_1=\dfrac{a_1q^2}{1-q}$ 约去 $a_1$ 得
\nopagebreak[4]
\[1=\frac{q^2}{1-q}\iff q^2+q-1=0,\]
\Needspace{5\baselineskip}
解得 $q=\dfrac{-1\pm\sqrt5}{2}$。结合 $|q|<1$ 得
\nopagebreak[4]
\[q=\frac{\sqrt5-1}{2}.\]

\gkanswer{答案：$\dfrac{\sqrt5-1}{2}$。由 $a_1=\dfrac{a_1q^2}{1-q}$ 得 $q^2+q-1=0$，结合 $|q|<1$ 取正根。}

\end{answer}
\begin{answer}
\gksolhead{9}{}

\textbf{题目}：

若 $f(x)=x^{\frac23}-x^{-\frac12}$，则满足 $f(x)<0$ 的 $x$ 的取值范围是\gkblank{}。

\Needspace{4\baselineskip}
\textbf{解析}

\Needspace{5\baselineskip}
由 $x^{-\frac12}$ 有意义得 $x>0$。不等式 $f(x)<0$ 即
\nopagebreak[4]
\[x^{\frac23}<x^{-\frac12},\]
\Needspace{5\baselineskip}
两边同乘 $x^{\frac12}>0$ 得
\nopagebreak[4]
\[x^{\frac23+\frac12}=x^{\frac76}<1,\]
故 $0<x<1$，取值范围为 $(0,1)$。

\gkanswer{答案：$(0,1)$。$x>0$，$x^{\frac23}<x^{-\frac12}\iff x^{\frac76}<1\iff0<x<1$。}

\end{answer}
\begin{answer}
\gksolhead{10}{}

\textbf{题目}：

为强化安全意识，某商场拟在未来的连续 $10$ 天中随机选择 $3$ 天进行紧急疏散演练，则选择的 $3$ 天恰好为连续 $3$ 天的概率是\gkblank{}（结果用最简分数表示）。

\Needspace{4\baselineskip}
\textbf{解析}

\Needspace{5\baselineskip}
从 $10$ 天中任选 $3$ 天，共有
\nopagebreak[4]
\[\mathrm C_{10}^3=\frac{10\times9\times8}{3\times2\times1}=120\]
种等可能的选法。

\Needspace{5\baselineskip}
$3$ 天恰好连续的情形为 $\{1,2,3\}$，$\{2,3,4\}$，$\cdots$，$\{8,9,10\}$，共 $8$ 种。故所求概率
\nopagebreak[4]
\[P=\frac8{120}=\frac1{15}.\]

\gkanswer{答案：$\dfrac1{15}$。连续三天共 $8$ 种，$\dfrac8{\mathrm C_{10}^3}=\dfrac8{120}=\dfrac1{15}$。}

\end{answer}
\begin{answer}
\gksolhead{11}{}

\textbf{题目}：

已知互异的复数 $a$，$b$ 满足 $ab\neq0$，集合 $\{a,b\}=\{a^2,b^2\}$，则 $a+b=$\gkblank{}。

\Needspace{4\baselineskip}
\textbf{解析}

由 $\{a,b\}=\{a^2,b^2\}$ 且 $a\neq b$，只有两种对应方式。

若 $a^2=a$ 且 $b^2=b$，则 $a$，$b\in\{0,1\}$，与 $ab\neq0$ 及 $a\neq b$ 矛盾。

\Needspace{5\baselineskip}
若 $a^2=b$ 且 $b^2=a$，则 $a^4=a$，即
\nopagebreak[4]
\[a(a-1)(a^2+a+1)=0 .\]
\Needspace{5\baselineskip}
由 $ab\neq0$ 得 $a\neq0$；由 $a\neq b$ 得 $a\neq1$（否则 $b=a^2=1=a$）。故 $a^2+a+1=0$，即 $a$ 是三次单位根 $\omega$（$\omega\neq1$，$\omega^3=1$），且 $b=a^2$。于是
\nopagebreak[4]
\[a+b=a+a^2=-1 .\]

\gkanswer{答案：$-1$。必有 $a^2=b$，$b^2=a$，得 $a^2+a+1=0$，故 $a+b=a+a^2=-1$。}

\end{answer}
\begin{answer}
\gksolhead{12}{}

\textbf{题目}：

设常数 $a$ 使方程 $\sin x+\sqrt3\cos x=a$ 在闭区间 $[0,2\pi]$ 上恰有三个解 $x_1$，$x_2$，$x_3$，则 $x_1+x_2+x_3=$\gkblank{}。

\Needspace{4\baselineskip}
\textbf{解析}

\Needspace{5\baselineskip}
由辅助角公式
\nopagebreak[4]
\[\sin x+\sqrt3\cos x=2\sin\left(x+\frac\pi3\right).\]
\Needspace{5\baselineskip}
令 $t=x+\dfrac\pi3$，则 $t\in\left[\dfrac\pi3,\dfrac{7\pi}3\right]$，原方程化为
\nopagebreak[4]
\[2\sin t=a .\]
$\sin t$ 在该区间上分三段单调：在 $\left[\dfrac\pi3,\dfrac\pi2\right]$ 上由 $\dfrac{\sqrt3}2$ 增到 $1$，在 $\left[\dfrac\pi2,\dfrac{3\pi}2\right]$ 上由 $1$ 减到 $-1$，在 $\left[\dfrac{3\pi}2,\dfrac{7\pi}3\right]$ 上由 $-1$ 增到 $\dfrac{\sqrt3}2$。逐段计数得：

\begin{itemize}
\item 当 $\dfrac a2=\dfrac{\sqrt3}2$ 时，三解为 $t=\dfrac\pi3$，$\dfrac{2\pi}3$，$\dfrac{7\pi}3$，共 $3$ 个；
\item 当 $\dfrac{\sqrt3}2<\dfrac a2\leqslant1$ 时，只有前两段各给出一解（$\dfrac a2=1$ 时两解重合），共 $2$ 个；
\item 当 $-1<\dfrac a2<\dfrac{\sqrt3}2$ 时，只有后两段各给出一解，共 $2$ 个；$\dfrac a2=-1$ 时仅一解。
\end{itemize}

\Needspace{5\baselineskip}
故恰有三个解当且仅当 $a=\sqrt3$，此时
\nopagebreak[4]
\[t_1=\frac\pi3,\qquad t_2=\frac{2\pi}3,\qquad t_3=\frac{7\pi}3,\]
\Needspace{5\baselineskip}
即 $x_1=0$，$x_2=\dfrac\pi3$，$x_3=2\pi$，从而
\nopagebreak[4]
\[x_1+x_2+x_3=0+\frac\pi3+2\pi=\frac{7\pi}3 .\]

\gkanswer{答案：$\dfrac{7\pi}3$。$2\sin\left(x+\dfrac\pi3\right)=a$ 恰三解当且仅当 $a=\sqrt3$，三解为 $0$，$\dfrac\pi3$，$2\pi$。}

\end{answer}
\begin{answer}
\gksolhead{13}{}

\textbf{题目}：

某游戏的得分为 $1,2,3,4,5$，随机变量 $\xi$ 表示小白玩该游戏的得分。若 $E(\xi)=4.2$，则小白得 $5$ 分的概率至少为\gkblank{}。

\Needspace{4\baselineskip}
\textbf{解析}

\Needspace{5\baselineskip}
记 $p_k=P(\xi=k)$（$k=1,2,3,4,5$），则 $\sum_{k=1}^5p_k=1$。由得分不超过 $5$、且只有得 $5$ 分时才可能超过 $4$ 分，
\nopagebreak[4]
\[4.2=E(\xi)=\sum_{k=1}^5kp_k\leqslant4\sum_{k=1}^4p_k+5p_5=4(1-p_5)+5p_5=4+p_5,\]
故 $p_5\geqslant0.2=\dfrac15$。

取 $p_4=0.8$，$p_5=0.2$，$p_1=p_2=p_3=0$，则 $E(\xi)=4\times0.8+5\times0.2=4.2$，下界可以取到。故小白得 $5$ 分的概率至少为 $\dfrac15$。

\gkanswer{答案：$\dfrac15$。$4.2=E(\xi)\leqslant4+p_5$，故 $p_5\geqslant0.2$，且 $p_5=0.2$，$p_4=0.8$ 时可达。}

\end{answer}
\begin{answer}
\gksolhead{14}{}

\textbf{题目}：

已知曲线 $C:x=-\sqrt{4-y^2}$，直线 $l:x=6$。若对于点 $A(m,0)$，存在 $C$ 上的点 $P$ 和 $l$ 上的 $Q$ 使得 $\overrightarrow{AP}+\overrightarrow{AQ}=\vec{0}$，则 $m$ 的取值范围为\gkblank{}。

\Needspace{4\baselineskip}
\textbf{解析}

\Needspace{5\baselineskip}
$C:x=-\sqrt{4-y^2}$ 是圆 $x^2+y^2=4$ 在 $y$ 轴左侧的半圆（含端点），故 $P(x_P,y_P)$ 满足
\nopagebreak[4]
\[x_P\in[-2,0],\qquad x_P^2+y_P^2=4 .\]

\Needspace{5\baselineskip}
设 $Q(6,y_Q)$。由 $\overrightarrow{AP}+\overrightarrow{AQ}=\vec{0}$ 知 $A$ 是线段 $PQ$ 的中点，故
\nopagebreak[4]
\[m=\frac{x_P+6}{2},\qquad 0=\frac{y_P+y_Q}{2}.\]
\Needspace{5\baselineskip}
由 $x_P\in[-2,0]$ 得
\nopagebreak[4]
\[m\in\left[\frac{-2+6}2,\frac{0+6}2\right]=[2,3].\]
反之，任给 $m\in[2,3]$，取 $x_P=2m-6\in[-2,0]$，令 $y_P=\sqrt{4-x_P^2}$，$y_Q=-y_P$，则 $P\in C$，$Q(6,y_Q)\in l$ 且 $\overrightarrow{AP}+\overrightarrow{AQ}=\vec{0}$。故 $m$ 的取值范围为 $[2,3]$。

\gkanswer{答案：$[2,3]$。$A$ 为 $PQ$ 中点，$m=\dfrac{x_P+6}2$，而 $x_P\in[-2,0]$。}

\end{answer}
\begin{answer}
\gksolhead{15}{}

\textbf{题目}：

设 $a$，$b\in\mathbf R$，则“$a+b>4$”是“$a>2$ 且 $b>2$”的\mbox{（　　）}

（A）充分非必要条件；（B）必要非充分条件；（C）充要条件；（D）既不充分也不必要条件。

\Needspace{4\baselineskip}
\textbf{解析}

若 $a>2$ 且 $b>2$，则 $a+b>4$，故“$a+b>4$”是“$a>2$ 且 $b>2$”的必要条件。

反之，取 $a=1$，$b=4$，则 $a+b=5>4$，但 $a>2$ 不成立，故“$a+b>4$”不是充分条件。

所以“$a+b>4$”是“$a>2$ 且 $b>2$”的必要非充分条件。

选 B。

\gkanswer{答案：B。$a>2$ 且 $b>2$ 推出 $a+b>4$，反之 $a=1$，$b=4$ 为反例。}

\end{answer}
\begin{answer}
\gksolhead{16}{}

\textbf{题目}：

如图，四个棱长为 $1$ 的正方体排成一个正四棱柱，$AB$ 是一条侧棱，$P_i$（$i=1,2,\cdots,8$）是上底面上其余的八个点，则 $\overrightarrow{AB}\cdot\overrightarrow{AP_i}$（$i=1,2,\cdots,8$）的不同值的个数为\mbox{（　　）}
\nopagebreak[4]
\begin{center}% ============================================================
% f2014_shanghai_science_q16_prism.tex
% 2014 上海卷（理科）第 16 题：四个棱长为 1 的正方体排成的正四棱柱
% 来源 PDF：普通高考/2014/2014上海理.pdf
% 原页：第 1 页（A3），印面 915，第 16 题题干右侧
% 底面 2×2、高 1；AB 为一条侧棱（A 在下、B 在上），
% 上底面除 B 外的八个格点为 P1—P8：
%   P1、P2 在左边（自 B 起），P3、P6 在前边，P4 为面心，
%   P5、P8 在后边，P7 在右边。
% 可见棱与上底面网格线为实线，被遮挡的棱与内部竖棱为虚线。
% 平行投影：前边方向取 1 单位 =（1.025,-0.15），竖直方向取 1 单位 =（0,-0.95）。
% ============================================================
\begin{tikzpicture}[>=Stealth, line width=0.65pt, font=\small, scale=1.1,
  line cap=round, line join=round]
  \coordinate (B) at (0,2.30);
  % 上底面九个格点
  \coordinate (P1) at ($(B)+(0.55,0.55)$);
  \coordinate (P2) at ($(B)+(1.10,1.10)$);
  \coordinate (P3) at ($(B)+(2.05,-0.30)$);
  \coordinate (P4) at ($(B)+(2.60,0.25)$);
  \coordinate (P5) at ($(B)+(3.15,0.80)$);
  \coordinate (P6) at ($(B)+(4.10,-0.60)$);
  \coordinate (P7) at ($(B)+(4.65,-0.05)$);
  \coordinate (P8) at ($(B)+(5.20,0.50)$);
  % 下底面九个格点（竖直下移 0.95）
  \coordinate (A)   at ($(B)+(0,-0.95)$);
  \coordinate (P1b) at ($(P1)+(0,-0.95)$);
  \coordinate (P2b) at ($(P2)+(0,-0.95)$);
  \coordinate (P3b) at ($(P3)+(0,-0.95)$);
  \coordinate (P4b) at ($(P4)+(0,-0.95)$);
  \coordinate (P5b) at ($(P5)+(0,-0.95)$);
  \coordinate (P6b) at ($(P6)+(0,-0.95)$);
  \coordinate (P7b) at ($(P7)+(0,-0.95)$);
  \coordinate (P8b) at ($(P8)+(0,-0.95)$);

  % 可见的外轮廓、可见棱与前、右底边
  \draw (A) -- (B) -- (P2) -- (P8) -- (P8b) -- (P6b) -- (A);
  \draw (B) -- (P6) -- (P8);
  \draw (P6) -- (P6b);
  \draw (P3) -- (P3b);
  \draw (P7) -- (P7b);

  % 上底面 2×2 网格线（过面心 P4）
  \draw (P3) -- (P5);
  \draw (P1) -- (P7);

  % 被遮挡的棱：后左竖棱、底面后边与左边、三条背面/内部竖棱
  \draw[dashed] (P2) -- (P2b);
  \draw[dashed] (P2b) -- (P8b);
  \draw[dashed] (P2b) -- (A);
  \draw[dashed] (P1) -- (P1b);
  \draw[dashed] (P4) -- (P4b);
  \draw[dashed] (P5) -- (P5b);

  % 顶点标注
  \fill (A) circle (0.85pt) node[below left=-1pt] {$A$};
  \fill (B) circle (0.85pt) node[left=-1pt] {$B$};
  \fill (P1) circle (0.85pt) node[left=-2pt] {$P_1$};
  \fill (P2) circle (0.85pt) node[above left=-1pt] {$P_2$};
  \fill (P3) circle (0.85pt) node[below left=-2pt] {$P_3$};
  \fill (P4) circle (0.85pt) node[above left=-2pt] {$P_4$};
  \fill (P5) circle (0.85pt) node[above=-1pt] {$P_5$};
  \fill (P6) circle (0.85pt) node[below left=-2pt] {$P_6$};
  \fill (P7) circle (0.85pt) node[right=-1pt] {$P_7$};
  \fill (P8) circle (0.85pt) node[above right=-1pt] {$P_8$};
\end{tikzpicture}
\end{center}
\nopagebreak[4]
（A）$1$；（B）$2$；（C）$4$；（D）$8$。

\Needspace{4\baselineskip}
\textbf{解析}

四个棱长为 $1$ 的正方体要排成上底面含九个格点的正四棱柱，只能是底面 $2\times2$、高 $1$ 的正四棱柱。

\Needspace{5\baselineskip}
以 $A$ 为原点，$\overrightarrow{AB}$ 的方向为 $z$ 轴正方向建立空间直角坐标系，则
\nopagebreak[4]
\[A(0,0,0),\qquad B(0,0,1),\]
\Needspace{5\baselineskip}
上底面（$z=1$）上的九个格点为 $(i,j,1)$（$i,j\in\{0,1,2\}$），其中 $B$ 对应 $(0,0,1)$，其余八个点 $P_i$ 的坐标 $(x_i,y_i,1)$ 恰好遍历另外八组 $(i,j)$。于是
\nopagebreak[4]
\[\overrightarrow{AB}=(0,0,1),\qquad \overrightarrow{AP_i}=(x_i,y_i,1),\]
\nopagebreak[4]
\[\overrightarrow{AB}\cdot\overrightarrow{AP_i}=0\cdot x_i+0\cdot y_i+1\cdot1=1 .\]
八个点给出同一个值，故不同值的个数为 $1$。

选 A。

\gkanswer{答案：A。取 $A(0,0,0)$，$B(0,0,1)$，八个 $P_i$ 均在 $z=1$ 上，故 $\overrightarrow{AB}\cdot\overrightarrow{AP_i}\equiv1$。}

\end{answer}
\begin{answer}
\gksolhead{17}{}

\textbf{题目}：

\Needspace{5\baselineskip}
已知 $P_1(a_1,b_1)$ 与 $P_2(a_2,b_2)$ 是直线 $y=kx+1$（$k$ 为常数）上两个不同的点，则关于 $x$ 和 $y$ 的方程组
\nopagebreak[4]
\[\begin{cases}
a_1x+b_1y=1,\\
a_2x+b_2y=1
\end{cases}\]
的解的情况是\mbox{（　　）}

（A）无论 $k$，$P_1$，$P_2$ 如何，总是无解；（B）无论 $k$，$P_1$，$P_2$ 如何，总有唯一解；（C）存在 $k$，$P_1$，$P_2$，使之恰有两解；（D）存在 $k$，$P_1$，$P_2$，使之有无穷多解。

\Needspace{4\baselineskip}
\textbf{解析}

\Needspace{5\baselineskip}
由 $P_i(a_i,b_i)$ 在直线 $y=kx+1$ 上得
\nopagebreak[4]
\[b_i=ka_i+1\qquad(i=1,2).\]
\Needspace{5\baselineskip}
方程组的系数行列式
\nopagebreak[4]
\[D=\begin{vmatrix}a_1&b_1\\a_2&b_2\end{vmatrix}=a_1b_2-a_2b_1=a_1(ka_2+1)-a_2(ka_1+1)=a_1-a_2 .\]
若 $a_1=a_2$，则 $b_1=ka_1+1=ka_2+1=b_2$，即 $P_1$，$P_2$ 重合，与题设矛盾。故 $D=a_1-a_2\neq0$，方程组恒有唯一解。

\Needspace{5\baselineskip}
事实上取 $x=-k$，$y=1$，则
\nopagebreak[4]
\[a_ix+b_iy=-ka_i+(ka_i+1)=1\qquad(i=1,2),\]
故解就是 $(-k,1)$，无论 $k$，$P_1$，$P_2$ 如何都是唯一解。

选 B。

\gkanswer{答案：B。由 $b_i=ka_i+1$ 得 $D=a_1-a_2\neq0$，且 $(-k,1)$ 恒为解。}

\end{answer}
\begin{answer}
\gksolhead{18}{}

\textbf{题目}：

\Needspace{5\baselineskip}
设
\nopagebreak[4]
\[f(x)=\begin{cases}
(x-a)^2, & x\leqslant0,\\
x+\dfrac1x+a, & x>0,
\end{cases}\]
若 $f(0)$ 是 $f(x)$ 的最小值，则 $a$ 的取值范围为\mbox{（　　）}

（A）$[-1,2]$；（B）$[-1,0]$；（C）$[1,2]$；（D）$[0,2]$。

\Needspace{4\baselineskip}
\textbf{解析}

$f(0)=a^2$。

\Needspace{5\baselineskip}
当 $x\leqslant0$ 时，$f(x)=(x-a)^2$。要求 $f(x)\geqslant f(0)=a^2$ 对一切 $x<0$ 成立：
\nopagebreak[4]
\[(x-a)^2\geqslant a^2\iff x^2-2ax\geqslant0\iff x(x-2a)\geqslant0 .\]
\Needspace{5\baselineskip}
因 $x<0$，该式等价于 $x-2a\leqslant0$，即 $a\geqslant\dfrac x2$ 对一切 $x<0$ 成立，故
\nopagebreak[4]
\[a\geqslant0 .\]

\Needspace{5\baselineskip}
当 $x>0$ 时，$f(x)=x+\dfrac1x+a\geqslant2+a$（$x=1$ 时取等号）。要求
\nopagebreak[4]
\[a^2\leqslant2+a\iff a^2-a-2\leqslant0\iff-1\leqslant a\leqslant2 .\]

两条件合并得 $0\leqslant a\leqslant2$。

选 D。

\gkanswer{答案：D。$x\leqslant0$ 一侧要求 $a\geqslant0$；$x>0$ 一侧要求 $a^2\leqslant a+2$，即 $-1\leqslant a\leqslant2$。}

\end{answer}
\begin{answer}
\gksolhead{19}{}

\textbf{题目}：

底面边长为 $2$ 的正三棱锥 $P\text{-}ABC$，其表面展开图是三角形 $P_1P_2P_3$，如图，求 $\triangle P_1P_2P_3$ 的各边长及此三棱锥的体积 $V$。
\nopagebreak[4]
\begin{center}% ============================================================
% f2014_shanghai_science_q19_net.tex
% 2014 上海卷（理科）第 19 题：正三棱锥 P-ABC 的表面展开图
% 来源 PDF：普通高考/2014/2014上海理.pdf
% 原页：第 1 页（A3），印面 915，第 19 题题干下方（右栏下部）
% 展开图为三角形 P1P2P3：A、C 分别为 P1P3、P3P2 的中点，
% B 为 P1P2 的中点；△ABC 为底面（实线），大三角形三边为
% 展开图外边界（虚线）。
% ============================================================
\begin{tikzpicture}[>=Stealth, line width=0.65pt, font=\small, scale=0.78,
  line cap=round, line join=round]
  \coordinate (P1) at (0,0);
  \coordinate (P2) at (4,0);
  \coordinate (P3) at (2,3.4641);
  \coordinate (A)  at (1,1.73205);
  \coordinate (B)  at (2,0);
  \coordinate (C)  at (3,1.73205);

  % 展开图外边界（虚线）
  \draw[dashed] (P1) -- (P2) -- (P3) -- cycle;

  % 底面三角形（折痕，实线）
  \draw (A) -- (B) -- (C) -- cycle;

  % 顶点标注
  \fill (P1) circle (0.85pt) node[below left=-1pt] {$P_1$};
  \fill (P2) circle (0.85pt) node[below right=-1pt] {$P_2$};
  \fill (P3) circle (0.85pt) node[above=-1pt] {$P_3$};
  \fill (A) circle (0.85pt) node[above left=-1pt] {$A$};
  \fill (B) circle (0.85pt) node[below=-1pt] {$B$};
  \fill (C) circle (0.85pt) node[above right=-1pt] {$C$};
\end{tikzpicture}
\end{center}
\nopagebreak[4]
\Needspace{4\baselineskip}
\textbf{解析}

设侧棱长为 $l$，侧面等腰三角形的底角为 $\theta$（侧面的底边就是底面边长 $2$）。

\Needspace{5\baselineskip}
展开图中，$A$，$B$，$C$ 分别落在 $P_1P_3$，$P_1P_2$，$P_2P_3$ 上，三个侧面分别以 $P_3$，$P_1$，$P_2$ 为顶点。展开图是三角形，说明 $B$ 处三个角拼成平角：
\nopagebreak[4]
\[\angle P_1BA+\angle ABC+\angle CBP_2=180^\circ .\]
\Needspace{5\baselineskip}
由正三棱锥的对称性三个侧面全等，$\angle P_1BA=\angle CBP_2=\theta$，而底面是正三角形，$\angle ABC=60^\circ$，故
\nopagebreak[4]
\[2\theta+60^\circ=180^\circ,\qquad\theta=60^\circ .\]
\Needspace{5\baselineskip}
于是每个侧面都是边长为 $2$ 的正三角形，$l=2$，从而
\nopagebreak[4]
\[P_1P_2=P_2P_3=P_3P_1=2l=4 .\]

\Needspace{5\baselineskip}
此时正三棱锥的各棱长均为 $2$，是正四面体。底面外接圆半径
\nopagebreak[4]
\[r=\frac{2}{2\sin60^\circ}=\frac{2}{\sqrt3},\]
\Needspace{5\baselineskip}
高
\nopagebreak[4]
\[h=\sqrt{l^2-r^2}=\sqrt{4-\frac43}=\frac{2\sqrt6}{3},\]
\Needspace{5\baselineskip}
故
\nopagebreak[4]
\[V=\frac13\cdot\frac{\sqrt3}{4}\cdot2^2\cdot\frac{2\sqrt6}{3}=\frac13\cdot\sqrt3\cdot\frac{2\sqrt6}{3}=\frac{2\sqrt2}{3} .\]

\gkanswer{答案：$P_1P_2=P_2P_3=P_3P_1=4$，$V=\dfrac{2\sqrt2}3$。展开图为三角形迫使侧面为正三角形，侧棱 $l=2$，棱长全为 $2$。}

\end{answer}
\begin{answer}
\gksolhead{20}{}

\textbf{题目}：

设常数 $a\geqslant0$，函数 $f(x)=\dfrac{2^x+a}{2^x-a}$。

（1） 若 $a=4$，求函数 $y=f(x)$ 的反函数 $y=f^{-1}(x)$；

（2） 根据 $a$ 的不同取值，讨论函数 $y=f(x)$ 的奇偶性，并说明理由。

\Needspace{4\baselineskip}
\textbf{解析}

\gksubpart{（1）}

\Needspace{5\baselineskip}
$a=4$ 时
\nopagebreak[4]
\[f(x)=\frac{2^x+4}{2^x-4},\]
\Needspace{5\baselineskip}
定义域为 $x\neq2$。设 $y=\dfrac{2^x+4}{2^x-4}$，则
\nopagebreak[4]
\[y(2^x-4)=2^x+4\iff (y-1)2^x=4(y+1)\iff2^x=\frac{4(y+1)}{y-1},\]
\Needspace{5\baselineskip}
故
\nopagebreak[4]
\[x=\log_2\frac{4(y+1)}{y-1},\]
\Needspace{5\baselineskip}
其中需 $y\neq1$ 且 $\dfrac{y+1}{y-1}>0$，即 $y<-1$ 或 $y>1$。所以
\nopagebreak[4]
\[f^{-1}(x)=\log_2\frac{4(x+1)}{x-1},\qquad x\in(-\infty,-1)\cup(1,+\infty).\]

\gksubpart{（2）}

\Needspace{5\baselineskip}
对定义域内的 $x$，
\nopagebreak[4]
\[f(-x)=\frac{2^{-x}+a}{2^{-x}-a}=\frac{1+a\cdot2^x}{1-a\cdot2^x}=-\frac{a\cdot2^x+1}{a\cdot2^x-1}.\]

\begin{itemize}
\item 当 $a=0$ 时，$f(x)=\dfrac{2^x}{2^x}=1$，定义域为 $\mathbf R$，$f(-x)=1=f(x)$，是偶函数；
\item 当 $a=1$ 时，定义域为 $\{x\mid x\neq0\}$，关于原点对称，且
\end{itemize}
\nopagebreak[4]
\[f(-x)=-\frac{2^x+1}{2^x-1}=-f(x),\]
是奇函数；
\begin{itemize}
\item 当 $a>0$ 且 $a\neq1$ 时，定义域为 $\{x\mid x\neq\log_2a\}$，因 $\log_2a\neq0$，定义域不关于原点对称，故既不是奇函数也不是偶函数。
\end{itemize}

\gkanswer{答案（1）：$f^{-1}(x)=\log_2\dfrac{4(x+1)}{x-1}$，定义域 $(-\infty,-1)\cup(1,+\infty)$。由 $2^x=\dfrac{4(y+1)}{y-1}$ 解出。}

\gkanswer{答案（2）：$a=0$ 时为偶函数；$a=1$ 时为奇函数；$a>0$ 且 $a\neq1$ 时既非奇函数也非偶函数。$a>0$ 且 $a\neq1$ 时定义域 $\{x\mid x\neq\log_2a\}$ 不关于原点对称。}

\end{answer}
\begin{answer}
\gksolhead{21}{}

\textbf{题目}：

如图，某公司要在 $A$、$B$ 两地连线上的定点 $C$ 处建造广告牌 $CD$，其中 $D$ 为顶端，$AC$ 长 $35$ 米，$CB$ 长 $80$ 米。设点 $A$、$B$ 在同一水平面上，从 $A$ 和 $B$ 看 $D$ 的仰角分别为 $\alpha$ 和 $\beta$。
\nopagebreak[4]
\begin{center}% ============================================================
% f2014_shanghai_science_q21_billboard.tex
% 2014 上海卷（理科）第 21 题：广告牌 CD 的仰角示意图
% 来源 PDF：普通高考/2014/2014上海理.pdf
% 原页：第 2 页（A3），印面 916，第 21 题题干下方（左栏）
% A、C、B 在同一水平线上，AC=35 米、CB=80 米，D 为广告牌顶端；
% 图中 CD 为铅垂方向（第 (1) 问的情形），α、β 分别为从 A、B
% 看 D 的仰角。示意图按 AC:CB:CD≈35:80:28 的比例绘制。
% ============================================================
\begin{tikzpicture}[>=Stealth, line width=0.65pt, font=\small, scale=1.0,
  line cap=round, line join=round]
  \coordinate (A) at (0,0);
  \coordinate (C) at (3.50,0);
  \coordinate (B) at (11.50,0);
  \coordinate (D) at (3.50,2.80);

  % 水平线与广告牌
  \draw (A) -- (B);
  \draw (D) -- (C);
  \draw (A) -- (D);
  \draw (B) -- (D);

  % 仰角 α、β 的角标
  \draw (1.00,0) arc (0:38.66:1.00);
  \draw ([shift=(160.71:1.00)]B) arc (160.71:180:1.00);

  % 顶点标注
  \fill (A) circle (0.85pt) node[below left=-1pt] {$A$};
  \fill (C) circle (0.85pt) node[below=-1pt] {$C$};
  \fill (B) circle (0.85pt) node[below right=-1pt] {$B$};
  \fill (D) circle (0.85pt) node[above=-1pt] {$D$};
  \node at (1.50,0.45) {$\alpha$};
  \node at (10.30,0.62) {$\beta$};
\end{tikzpicture}
\end{center}
\nopagebreak[4]
（1） 设计中 $CD$ 是铅垂方向。若要求 $\alpha\geqslant2\beta$，问 $CD$ 的长至多为多少（结果精确到 $0.01$ 米）？

（2） 施工完成后，$CD$ 与铅垂方向有偏差。现在实测得 $\alpha=38.12^\circ$，$\beta=18.45^\circ$，求 $CD$ 的长（结果精确到 $0.01$ 米）。

\Needspace{4\baselineskip}
\textbf{解析}

\gksubpart{（1）}

\Needspace{5\baselineskip}
设 $CD=h$（米）。由 $CD$ 铅垂、$A$，$C$，$B$ 在同一水平线上，得
\nopagebreak[4]
\[\tan\alpha=\frac h{35},\qquad\tan\beta=\frac h{80}.\]
\Needspace{5\baselineskip}
由 $\alpha\geqslant2\beta$ 得 $\tan\alpha\geqslant\tan2\beta$（若 $2\beta\geqslant\dfrac\pi2$，则已有 $\alpha<\dfrac\pi2\leqslant2\beta$ 不满足条件，对应 $h\geqslant80>20\sqrt2$，不影响所得界），而
\nopagebreak[4]
\[\tan2\beta=\frac{2\tan\beta}{1-\tan^2\beta}=\frac{2\cdot\frac h{80}}{1-\left(\frac h{80}\right)^2}=\frac{160h}{6400-h^2},\]
\Needspace{5\baselineskip}
故
\nopagebreak[4]
\[\frac h{35}\geqslant\frac{160h}{6400-h^2}.\]
\Needspace{5\baselineskip}
因 $h>0$，约去 $h$ 得
\nopagebreak[4]
\[\frac1{35}\geqslant\frac{160}{6400-h^2}\iff6400-h^2\geqslant5600\iff h^2\leqslant800,\]
故 $h\leqslant20\sqrt2\approx28.28$。所以 $CD$ 的长至多为 $20\sqrt2$ 米，约 $28.28$ 米。

\gksubpart{（2）}

\Needspace{5\baselineskip}
设 $D$ 在地面上的投影为 $D'$，$AD'=x$ 米，$DD'=h$ 米，则
\nopagebreak[4]
\[\tan\alpha=\frac hx,\qquad\tan\beta=\frac h{115-x},\qquad x\in(0,115).\]
\Needspace{5\baselineskip}
由 $h=x\tan\alpha=(115-x)\tan\beta$ 得
\nopagebreak[4]
\[x=\frac{115\tan\beta}{\tan\alpha+\tan\beta}.\]
\Needspace{5\baselineskip}
取 $\tan38.12^\circ\approx0.7847$，$\tan18.45^\circ\approx0.3336$，则
\nopagebreak[4]
\[x\approx\frac{115\times0.3336}{0.7847+0.3336}\approx34.31,\qquad h=x\tan\alpha\approx34.31\times0.7847\approx26.92 .\]
\Needspace{5\baselineskip}
于是
\nopagebreak[4]
\[CD=\sqrt{(x-35)^2+h^2}\approx\sqrt{(-0.69)^2+26.92^2}\approx26.93 .\]
故施工完成后 $CD$ 的长约为 $26.93$ 米。

\gkanswer{答案（1）：至多 $20\sqrt2$ 米，约 $28.28$ 米。由 $\tan\alpha\geqslant\tan2\beta$ 得 $6400-h^2\geqslant5600$。}

\gkanswer{答案（2）：约 $26.93$ 米。由 $x=\dfrac{115\tan\beta}{\tan\alpha+\tan\beta}$、$h=x\tan\alpha$ 得水平偏差约 $0.69$ 米、高约 $26.92$ 米，$CD=\sqrt{0.69^2+26.92^2}$。}

\end{answer}
\begin{answer}
\gksolhead{22}{}

\textbf{题目}：

在平面直角坐标系 $xOy$ 中，对于直线 $l:ax+by+c=0$ 和点 $P_1(x_1,y_1)$，$P_2(x_2,y_2)$，记 $\eta=(ax_1+by_1+c)(ax_2+by_2+c)$。若 $\eta<0$，则称点 $P_1$，$P_2$ 被直线 $l$ 分隔。若曲线 $C$ 与直线 $l$ 没有公共点，且曲线 $C$ 上存在点 $P_1$，$P_2$ 被直线 $l$ 分隔，则称直线 $l$ 为曲线 $C$ 的一条分隔线。

（1） 求证：点 $A(1,2)$，$B(-1,0)$ 被直线 $x+y-1=0$ 分隔；

（2） 若直线 $y=kx$ 是曲线 $x^2-4y^2=1$ 的分隔线，求实数 $k$ 的取值范围；

（3） 动点 $M$ 到点 $Q(0,2)$ 的距离与到 $y$ 轴的距离之积为 $1$，设点 $M$ 的轨迹为 $E$，求证：通过原点的直线中，有且仅有一条直线是 $E$ 的分隔线。

\Needspace{4\baselineskip}
\textbf{解析}

\gksubpart{（1）}

\Needspace{5\baselineskip}
对直线 $x+y-1=0$ 有 $a=b=1$，$c=-1$，于是
\nopagebreak[4]
\[\eta=(1+2-1)\times(-1+0-1)=2\times(-2)=-4<0,\]
故点 $A(1,2)$，$B(-1,0)$ 被直线 $x+y-1=0$ 分隔。

\gksubpart{（2）}

\Needspace{5\baselineskip}
把 $y=kx$ 代入 $x^2-4y^2=1$ 得
\nopagebreak[4]
\[(1-4k^2)x^2=1 .\]
当 $1-4k^2>0$ 即 $|k|<\dfrac12$ 时，$x=\pm\dfrac1{\sqrt{1-4k^2}}$，直线与曲线有公共点；当 $|k|\geqslant\dfrac12$ 时方程无解，直线与曲线没有公共点。

\Needspace{5\baselineskip}
再验证 $|k|\geqslant\dfrac12$ 时曲线上的点确能被分隔。取 $t>1$，令 $u=\dfrac{\sqrt{t^2-1}}2$，则
\nopagebreak[4]
\[P_1(t,u),\qquad P_2(-t,u)\]
\Needspace{5\baselineskip}
都在 $x^2-4y^2=1$ 上，且
\nopagebreak[4]
\[\eta=(kt-u)(-kt-u)=u^2-k^2t^2 .\]
由 $|k|\geqslant\dfrac12$ 得 $k^2t^2\geqslant\dfrac{t^2}4>\dfrac{t^2-1}4=u^2$，故 $\eta<0$，即 $P_1$，$P_2$ 被直线 $y=kx$ 分隔。

所以 $k$ 的取值范围为 $\left(-\infty,-\dfrac12\right]\cup\left[\dfrac12,+\infty\right)$。

\gksubpart{（3）}

\Needspace{5\baselineskip}
由题意，$E$ 的方程为
\nopagebreak[4]
\[\sqrt{x^2+(y-2)^2}\cdot|x|=1\qquad(x\neq0),\]
\Needspace{5\baselineskip}
即
\nopagebreak[4]
\[x^2\left[x^2+(y-2)^2\right]=1 .\]
$E$ 关于 $y$ 轴对称，且 $(1,2)$，$(-1,2)$ 都在 $E$ 上。

\Needspace{5\baselineskip}
先看 $y$ 轴（即 $x=0$）。$E$ 上的点 $x\neq0$，故 $y$ 轴与 $E$ 没有公共点；而对直线 $x=0$，取 $P_1(1,2)$，$P_2(-1,2)\in E$，有
\nopagebreak[4]
\[\eta=1\times(-1)=-1<0,\]
故 $y$ 轴是 $E$ 的一条分隔线。

\Needspace{5\baselineskip}
再看任意一条不是 $y$ 轴的过原点的直线 $y=kx$。对 $x>0$ 令
\nopagebreak[4]
\[F(x)=x^2\left[x^2+(kx-2)^2\right],\]
\Needspace{5\baselineskip}
则 $F$ 在 $[0,1]$ 上连续，且
\nopagebreak[4]
\[F(0)=0<1,\qquad F(1)=1+(k-2)^2\geqslant1 .\]
由介值定理，存在 $x_0\in(0,1]$ 使 $F(x_0)=1$，即点 $(x_0,kx_0)$ 同时在 $E$ 与直线 $y=kx$ 上。所以这条直线与 $E$ 有公共点，不是分隔线。

综上，通过原点的直线中，有且仅有一条（即 $y$ 轴）是 $E$ 的分隔线。

\gkanswer{答案（1）：证明。$\eta=(1+2-1)(-1+0-1)=2\times(-2)=-4<0$，满足分隔定义。}

\gkanswer{答案（2）：$k\in\left(-\infty,-\dfrac12\right]\cup\left[\dfrac12,+\infty\right)$。$|k|<\dfrac12$ 时直线与双曲线相交；$|k|\geqslant\dfrac12$ 时无公共点，且取关于 $y$ 轴对称的两点可使 $\eta=u^2-k^2t^2<0$。}

\gkanswer{答案（3）：证明。$y$ 轴与 $E$ 无公共点且分隔 $(1,2)$ 与 $(-1,2)$；其余过原点的直线 $y=kx$ 均与 $E$ 相交，因为 $F(0)=0<1\leqslant F(1)$。}

\end{answer}
\begin{answer}
\gksolhead{23}{}

\textbf{题目}：

已知数列 $\{a_n\}$ 满足 $\dfrac13a_n\leqslant a_{n+1}\leqslant3a_n$，$n\in\mathbf N^*$，$a_1=1$。

（1） 若 $a_2=2$，$a_3=x$，$a_4=9$，求 $x$ 的取值范围；

（2） 设 $\{a_n\}$ 是公比为 $q$ 的等比数列，$S_n=a_1+a_2+\cdots+a_n$。若 $\dfrac13S_n\leqslant S_{n+1}\leqslant3S_n$，$n\in\mathbf N^*$，求 $q$ 的取值范围；

（3） 若 $a_1,a_2,\cdots,a_k$ 成等差数列，且 $a_1+a_2+\cdots+a_k=1000$，求正整数 $k$ 的最大值，以及 $k$ 取最大值时相应数列 $a_1,a_2,\cdots,a_k$ 的公差。

\Needspace{4\baselineskip}
\textbf{解析}

\gksubpart{（1）}

\Needspace{5\baselineskip}
由 $\dfrac13a_2\leqslant a_3\leqslant3a_2$ 得
\nopagebreak[4]
\[\frac23\leqslant x\leqslant6 .\]
\Needspace{5\baselineskip}
由 $\dfrac13a_3\leqslant a_4\leqslant3a_3$ 得
\nopagebreak[4]
\[\frac x3\leqslant9\leqslant3x\iff3\leqslant x\leqslant27 .\]
取交集得 $3\leqslant x\leqslant6$，即 $x$ 的取值范围为 $[3,6]$。

\gksubpart{（2）}

设 $a_n=q^{n-1}$（等比数列各项非零，$q\neq0$）。

\Needspace{5\baselineskip}
先看数列自身的条件。若 $q<0$，取 $n=2$ 得 $\dfrac13q\leqslant q^2\leqslant3q$，而 $q^2>0>3q$ 与右半式矛盾，故 $q>0$。此时约去 $q^{n-1}>0$ 得
\nopagebreak[4]
\[\frac13\leqslant q\leqslant3 .\]

\Needspace{5\baselineskip}
再看 $S_n$ 的条件。由 $S_{n+1}=S_n+q^n$ 及 $q>0$ 知 $S_n>0$，故 $\dfrac13S_n\leqslant S_{n+1}$ 自动成立；另一侧化为
\nopagebreak[4]
\[S_n+q^n\leqslant3S_n\iff q^n\leqslant2S_n .\]

\begin{itemize}
\item 若 $q=1$，则 $S_n=n$，$q^n=1\leqslant2n$ 对一切 $n\in\mathbf N^*$ 成立；
\item 若 $1<q\leqslant3$，则 $S_n=\dfrac{q^n-1}{q-1}$，条件化为
\end{itemize}
\nopagebreak[4]
\[q^n(q-1)\leqslant2(q^n-1)\iff q^n(3-q)\geqslant2 .\]
左端关于 $n$ 递增，故只需 $n=1$：$q(3-q)\geqslant2$，即 $(q-1)(q-2)\leqslant0$，得 $1<q\leqslant2$；
\begin{itemize}
\item 若 $\dfrac13\leqslant q<1$，则 $S_n=\dfrac{1-q^n}{1-q}$，条件化为
\end{itemize}
\nopagebreak[4]
\[q^n(1-q)\leqslant2(1-q^n)\iff q^n(3-q)\leqslant2 .\]
左端关于 $n$ 递减，故只需 $n=1$：$q(3-q)\leqslant2$，即 $q\leqslant1$ 或 $q\geqslant2$，在 $\left[\dfrac13,1\right)$ 上恒成立。

综上，$q$ 的取值范围为 $\left[\dfrac13,2\right]$。

\gksubpart{（3）}

\Needspace{5\baselineskip}
设公差为 $d$，则 $a_j=1+(j-1)d$（$j=1,2,\cdots,k$），
\nopagebreak[4]
\[1000=\frac{k\left[1+\left(1+(k-1)d\right)\right]}2=\frac{k\left[2+(k-1)d\right]}2,\]
\Needspace{5\baselineskip}
故
\nopagebreak[4]
\[d=\frac{2(1000-k)}{k(k-1)} .\]
\Needspace{5\baselineskip}
对 $j=1,\cdots,k-1$ 需
\nopagebreak[4]
\[\frac13a_j\leqslant a_j+d\leqslant3a_j\iff-\frac23a_j\leqslant d\leqslant2a_j .\]

\Needspace{5\baselineskip}
若 $d\geqslant0$，则 $a_j$ 递增，$d\leqslant2a_j$ 只需 $d\leqslant2a_1=2$，下界自动成立。由 $d\leqslant2$ 得
\nopagebreak[4]
\[1000-k\leqslant k(k-1)\iff k^2\geqslant1000,\]
结合 $d\geqslant0$ 即 $k\leqslant1000$，此时最大只能取到 $k=1000$。

\Needspace{5\baselineskip}
若 $d<0$，则 $k>1000$，$a_j$ 递减，上界 $d\leqslant2a_j$ 自动成立，下界 $d\geqslant-\dfrac23a_j$ 只需对最小的 $a_{k-1}=1+(k-2)d$ 成立：
\nopagebreak[4]
\[d\geqslant-\frac23\left[1+(k-2)d\right]\iff d(2k-1)\geqslant-2\iff d\geqslant-\frac2{2k-1} .\]
\Needspace{5\baselineskip}
代入 $d=\dfrac{2(1000-k)}{k(k-1)}<0$：
\nopagebreak[4]
\[\frac{2(1000-k)}{k(k-1)}\geqslant-\frac2{2k-1}\iff(1000-k)(2k-1)\geqslant-k(k-1)\]
\nopagebreak[4]
\[\iff-k^2+2000k-1000\geqslant0\iff k^2-2000k+1000\leqslant0,\]
\Needspace{5\baselineskip}
解得
\nopagebreak[4]
\[1000-\sqrt{999000}\leqslant k\leqslant1000+\sqrt{999000},\]
故整数 $k\leqslant1999$（因 $1000+\sqrt{999000}\approx1999.5$）。

\Needspace{5\baselineskip}
当 $k=1999$ 时，
\nopagebreak[4]
\[d=\frac{2(1000-1999)}{1999\times1998}=\frac{-1998}{1999\times1998}=-\frac1{1999},\]
\Needspace{5\baselineskip}
此时 $a_{1999}=1+1998\times\left(-\dfrac1{1999}\right)=\dfrac1{1999}>0$，且
\nopagebreak[4]
\[d=-\frac1{1999}\geqslant-\frac2{2\times1999-1}=-\frac2{3997},\]
各项条件均满足（$k$ 项之后的项取 $a_{n+1}=a_n$ 即可延续）。

所以正整数 $k$ 的最大值为 $1999$，此时相应数列的公差为 $d=-\dfrac1{1999}$。

\gkanswer{答案（1）：$[3,6]$。由 $\dfrac13x\leqslant9\leqslant3x$ 得 $3\leqslant x\leqslant27$，与 $\dfrac23\leqslant x\leqslant6$ 取交。}

\gkanswer{答案（2）：$\left[\dfrac13,2\right]$。数列条件给出 $\dfrac13\leqslant q\leqslant3$，和的条件化为 $q^n\leqslant2S_n$，分 $q=1$、$q>1$、$q<1$ 讨论。}

\gkanswer{答案（3）：$k$ 的最大值为 $1999$，此时公差 $d=-\dfrac1{1999}$。$d\geqslant0$ 时 $k\leqslant1000$；$d<0$ 时由 $d\geqslant-\dfrac2{2k-1}$ 得 $k^2-2000k+1000\leqslant0$。}

\end{answer}